% Section 6: Discussion (~1 page, ~900 words)
% Outline only — prose pass deferred until after Friday meeting.
% The "Avoiding the Curse of Diversity" paragraph is kept as concrete
% prose because the factual claim (independent-buffer construction =
% immune by design) is already locked in regardless of how Section 1
% framing shifts.

\section{Discussion}
\label{sec:discussion}

\subsection{Why the ensemble works in this regime}
\label{sec:d_why_works}

\guideblock{One-paragraph mechanism story. Each member's $Q$-table is
trained from a different stochastic trajectory through the same MDP,
so per-state action-value estimates have errors that are largely
uncorrelated across members. Voting averages those errors out; the
expected vote argmax converges toward the true argmax faster than any
single member. Cross-reference the per-member disagreement diagnostic
(Figure~\ref{fig:diag}, right) to show the diversity assumption is
satisfied in practice, not just in theory.}

\subsection{Why ensembling beats more single-agent compute}
\label{sec:d_why_beats}

\guideblock{One-paragraph paragraph framing the
compute-matched comparison. \texttt{baseline\_5m} plateaus around $54\%$
after $5{\times}10^6$ battles; ten agents at $5{\times}10^5$ battles
each cover the state space more broadly than one agent at $5{\times}10^6$
does, even though total simulated battles are identical at
$5{\times}10^6$. Tie this to the $K$-saturation curve
(Figure~\ref{fig:ksat}): the lift comes from \emph{state-space coverage
breadth}, not from extra wall-clock work. This is the
``coverage-versus-depth'' trade-off the paper makes explicit.}

\subsection{Avoiding the Curse of Diversity}
\label{sec:d_curse}

The recently-documented ``Curse of Diversity'' failure mode for
ensemble-based exploration \citep{lin2024curse} arises when ensemble
members share a replay buffer: members trained on largely off-policy
data relative to their own policy underperform a single-agent baseline.
By construction, the ensemble in this paper does not share a buffer or
any training data across members: each member is a separate process with
its own random seed, its own opponent battles, and its own $Q$-table.
The failure mode therefore does not apply, and our per-member solo win
rates (Figure~\ref{fig:diag}, left) do not underperform a single-agent
baseline.

\subsection{Limitations}
\label{sec:d_limitations}

\guideblock{Enumerate the honest gaps. Treat each as a real limitation,
not a teaser for hypothetical future work:
(1) Single opponent. All evaluations are against
\texttt{SimpleHeuristicsPlayer} in \texttt{gen4ou}; generalization across
opponents (e.g., \texttt{MaxBasePower} or a stronger learned agent) is
untested.
(2) Single 20-Pokemon pool. We do not study generalization across pools.
(3) Tabular only. The inference rule itself does not require a tabular
$Q$-function, but whether the $+12.6$ pp lift survives under deep value
approximation is an open empirical question.
(4) $K$ capped at 10 at evaluation. This is a workstation memory limit
($\approx 5.5$ GiB per member, ten members concurrent), not a
methodological choice. The $K$-saturation curve already plateaus by
$K{=}10$, so the saturating shape itself is robust to this cap, but a
tighter saturation point at higher $K$ cannot be ruled out from these
data.
(5) Homogeneous per-member configuration. Every member uses the fixed
winning configuration of the prior factorial study
(Appendix~\ref{app:inherited_components}); we do not vary per-member
hyperparameters and so cannot speak to whether the ensemble gain depends
on homogeneity.}
